
\documentclass[a4paper,11pt]{article}

\usepackage[utf8]{inputenc}
\usepackage[T1]{fontenc}
\usepackage[english]{babel}
\usepackage{times,lmodern,charter}
\usepackage{schooldocs}
\usepackage{enumerate}
\usepackage{multicol}
\usepackage{centeredline}
\usepackage{hyperref}
\usepackage{lipsum}

\definecolor{darkbrown}{rgb}{0.5,0.1,0.1}

\pagestyle{collection}
\title{The schooldocs package}
\subject{Example file for titles and styles}
\author{Antoine Missier}
\institute{Royal College of Pataphysics}
\date{December 28, 2023}
\renewcommand{\titlestyle}{\Huge} %\sffamily
\hypersetup{colorlinks, linkcolor=blue}
%\setcounter{tocdepth}{0}

\begin{document}

\maketitle

On the following pages we showcase various title formatting options
and the different styles. 
This initial page has been created using the \textbf{collection} style, for which
the main title style has been redefined through
\verb|\renewcommand{\titlestyle}{\Huge}|.


\medskip
In most examples (dated may 4, 2020), we have utilized
the \textsf{lmodern} package font, which pairs well with a sans-serif title style.
To achieve this, simply redefine the title style using
\smallskip
\centeredline{\verb|\renewcommand\titelstyle{\LARGE\bfseries\sffamily}|.}

\smallskip
For examples dated january 25 or december 28, 2023, we have opted 
for the \textsf{charter}, or	\textsf{times} fonts. 
For these fonts, the default title style setting 
\verb|\LARGE\bfseries| is suitable.


\seprule

\begin{center} \Large \color{darkbrown}
\bfseries Table of contents
\end{center}

\bigskip

\begin{multicols}{2}
\begin{enumerate}[\bfseries 1.]
	\item \hyperlink{schooldocs-example-center}{\bfseries The standard title}
	\item \hyperlink{schooldocs-example-left}{\bfseries A left-aligned title}
	\item \hyperlink{schooldocs-example-right}{\bfseries A right-aligned title}
	\item \hyperlink{schooldocs-example-norule}{\bfseries
		The standard title without rule}
	\item \hyperlink{schooldocs-example-boxed}{\bfseries A title in boxed shape}
	\item \hyperlink{schooldocs-example-seprule}{\bfseries
		The \texttt{\textbackslash seprule} and \texttt{\textbackslash correct} macros}
	\item \hyperlink{schooldocs-example-small}{\bfseries A ``small'' title}
	\item \hyperlink{schooldocs-example-classic}{\bfseries The classic style layout}
	\item \hyperlink{schooldocs-example-elegant}{\bfseries The elegant style layout}
	\item \hyperlink{schooldocs-example-modern}{\bfseries The modern style layout}
	\item \hyperlink{schooldocs-example-soft}{\bfseries The soft style layout}
	\item \hyperlink{schooldocs-example-exam}{\bfseries The exam style layout}
	\item \hyperlink{schooldocs-example-examcorrect}{\bfseries
		Exam style with the \texttt{\textbackslash correct} macro}
	\item \hyperlink{schooldocs-example-collection}{\bfseries The collection style layout}
	\item \hyperlink{schooldocs-example-identity}{\bfseries 
		The identity style in collection}
	\item \hyperlink{schooldocs-example-small-identity}{\bfseries 
		A ``small'' title with identity}
	\item \hyperlink{schooldocs-example-modern-small-identity}{\bfseries 
		Small title, identity and modern style}
	\item \hyperlink{schooldocs-example-left-identity}{\bfseries 
		A left-aligned title with identity}
\end{enumerate}
\end{multicols}

\newpage
\hypertarget{schooldocs-example-center}{}
\addcontentsline{toc}{section}{The standard title}

\fontfamily{lmr}\selectfont
\pagestyle{empty}
\lfoot{}
\cfoot{}
\rfoot{}
\renewcommand{\subjectstyle}{}
\renewcommand{\authorstyle}{}
\renewcommand{\datestyle}{}
\renewcommand{\titlestyle}{\LARGE\bfseries\sffamily}
\renewcommand\subjectstyle{\large}
\setlength{\titletopskip}{-1.32cm}
\setlength\titlesep{2\medskipamount}
\schooldocstitles
\title{The standard title}
\subject[Pathography / M2A]{Pathography / Master 2A}
\date{May 4, 2020}

\maketitle
\lipsum[1]

\newpage
\hypertarget{schooldocs-example-left}{}
\addcontentsline{toc}{section}{A left-aligned title}

\renewcommand{\titleflush}{flushleft}
\title{A left-aligned title}
\maketitle
\lipsum[1]

\newpage
\hypertarget{schooldocs-example-right}{}

\renewcommand{\titleflush}{flushright}
\title{A right-aligned title}
\maketitle
\lipsum[1]

\newpage
\hypertarget{schooldocs-example-norule}{}

\renewcommand{\titleflush}{center}
\title{The standard title without rule}
\maketitle[0pt]
\lipsum[1]

\newpage
\hypertarget{schooldocs-example-boxed}{}

\renewcommand{\titlestyle}{\LARGE\sffamily\bfseries\boxedshape}
\title{A title in boxed shape}
\maketitle
\lipsum[1]

\newpage
\hypertarget{schooldocs-example-seprule}{}

\renewcommand{\titlestyle}{\LARGE\bfseries\sffamily}
\title{The \textbackslash seprule and \textbackslash correct macros}
\correct
\maketitle
\lipsum[1]\seprule\lipsum[2]\seprule\lipsum[3]

\newpage
\hypertarget{schooldocs-example-small}{}

\fontfamily{bch}\selectfont
\title{A ``small'' title}
\date{January 25, 2023}
\makesmalltitle
\lipsum[1]

\newpage
\hypertarget{schooldocs-example-classic}{}

\fontfamily{lmr}\selectfont
\setcounter{page}{1}
\pagestyle{classic}
\title[The classic style]{The classic style layout}
\date{May 4, 2020}
\maketitle
\section*{Pellentesque a nulla}
\lipsum[1-3]
\section*{Vestibulum lectus}
\lipsum[5-9]

\newpage
\hypertarget{schooldocs-example-elegant}{}

\setcounter{page}{1}
\setcounter{section}{0}
\addtolength{\headheight}{-\baselineskip}
\pagestyle{elegant}
\title[The elegant style]{The elegant style layout}
\schooldocstitles
\maketitle
\section*{Pellentesque a nulla}
\lipsum[1-3]
\section*{Vestibulum lectus}
\lipsum[5-9]

\newpage
\hypertarget{schooldocs-example-modern}{}

\setcounter{page}{1}
\setcounter{section}{0}
\renewcommand{\headrulewidth}{0.4pt}
\pagestyle{modern}
\title[The modern style]{The modern style layout}
\maketitle
\section*{Pellentesque a nulla}
\lipsum[1-3]
\section*{Vestibulum lectus}
\lipsum[5-9]

\newpage
\hypertarget{schooldocs-example-soft}{}

\setcounter{page}{1}
\setcounter{section}{0}
\pagestyle{soft}
\title[The soft style]{The soft style layout}
\maketitle
\section*{Pellentesque a nulla}
\lipsum[1-3]
\section*{Vestibulum lectus}
\lipsum[5-9]

\newpage
\hypertarget{schooldocs-example-exam}{}

\setcounter{page}{1}
\setcounter{section}{0}
\fontfamily{ptm}\selectfont 
\pagestyle{exam}
\title[The exam style]{The exam style layout}
\date{December 28, 2023}
\subtitle{Duration of the test: 2\,h}
\maketitle
\section*{Pellentesque a nulla}
\lipsum[1-3]
\section*{Vestibulum lectus}
\lipsum[5-9]

\newpage
\hypertarget{schooldocs-example-examcorrect}{}

\setcounter{page}{1}
\setcounter{section}{0}
\correct
\maketitle
\section*{Pellentesque a nulla}
\lipsum[1-3]
\section*{Vestibulum lectus}
\lipsum[5-9]

\newpage
\hypertarget{schooldocs-example-collection}{}

\setcounter{page}{1}
\setcounter{section}{0}
\fontfamily{bch}\selectfont
\pagestyle{collection}
\title[The collection style]{The collection style layout}
\date{December 28, 2023}
\subtitle{Duration of the test: 2\,h}
\maketitle
\section*{Ultrices bibendum}
\lipsum
\lipsum[1-3]
\section*{Etiam facilisis}
\lipsum[1-5]

\newpage
\hypertarget{schooldocs-example-identity}{}

\thispagestyle{identity}
Here we present a single page with the \textbf{identity} style
within a document using the collection style.
\par\medskip
\textbf{Question:} Analyze the previous text and provide a written summary of its content.

\newpage
\lfoot{}
\cfoot{\thepage}
\rfoot{}
\renewcommand{\subjectstyle}{}
\renewcommand{\authorstyle}{}
\renewcommand{\datestyle}{}
\hypertarget{schooldocs-example-small-identity}{}

\setcounter{page}{1}
\renewcommand{\titlestyle}{\LARGE\bfseries}
\renewcommand\subjectstyle{\large}
\renewcommand{\headrulewidth}{0pt}
\setlength{\titletopskip}{-1.32cm}
\setlength\titlesep{2\medskipamount}
\title{A ``small'' title}
\date{January 25, 2023}
\thispagestyle{identity}
\schooldocstitles
\makesmalltitle
\lipsum[1]

\newpage
\hypertarget{schooldocs-example-modern-small-identity}{}

\setcounter{page}{1}
\setcounter{section}{0}
\title[The modern style]{The modern style layout}
\pagestyle{modern}
\thispagestyle{identity}
\makesmalltitle
\section*{Pellentesque a nulla}
\lipsum[1-3]
\section*{Vestibulum lectus}
\lipsum[5-6]

\newpage
\hypertarget{schooldocs-example-left-identity}{}

\setcounter{page}{1}
\lfoot{}
\cfoot{}
\rfoot{}
\renewcommand{\titleflush}{flushleft}
\title{A left-aligned title}
\pagestyle{empty}
\thispagestyle{identity}
\schooldocstitles
\maketitle
\lipsum[1]

\setcounter{page}{2} % for the lastpage field in headers or footers
\end{document}
